\section{Plans, forecasts and repeated errors}
\label{sec:temporal}

\paragraph{Plans need an opportunity before they can fail.} A fixed sample of 240 decisions, three per player per game drawn by hash rank from decisions not flagged in an earlier public-text extraction, was coded by one LLM reader for forward-looking statements in public messages and applied notes. Fifty-three sampled private fields are empty, all Opus's. The reader found forward-looking own-action passages in 59 records, statements about opponents' capacity or likely outcomes in 36, and references to earlier corrections in 25. Eight records state an explicit policy of buying named properties on landing. Six of these never meet a later purchase opportunity for the named deed. One is executed: in G2 Grok buys Mediterranean at its d256 opportunity. One fails at a real opportunity. At G16 d103 Grok writes ``buy Virginia on sight, then house pink''; at d159 it has \$371, Virginia costs \$160 and the menu offers the purchase, and its note says ``Buy at list is the only way to close pink'', yet the applied action starts an auction, which Gemini wins at \$380. Because the action contradicts its own note, this is not a deliberate change of plan. Plans also succeed over many turns: at G9 d64 Sol targets three houses on each dark blue as funds arrive and reaches that state at d163, thirty-one turns later. Appendix~\ref{app:temporal} gives these chronologies.

\paragraph{Most forecasts concern public capacity.} The opponent statements in the sample, and those in the incidents above, mostly compute what a rival can afford or would gain from a trade, using cash and holdings that every player can see. Some are right: at G2 d161 a forecast that withdrawing from an auction leaves New York to a named player is realized. Others have no test, as when Grok predicts at G19 d220 that rivals will not pay it to survive and no offer arrives before its elimination. These are financial deductions. The records contain no probability forecasts and no pre-registered targets, so they cannot score calibration or inference about hidden beliefs.

\paragraph{Corrections are local.} Among 114 decisions with an uneven-building error, 54 are followed within the same decision by a corrective request that contains the exact validator error. Fifty of these error decisions are followed by another build submission from the same player in the same game, and nine of those submissions repeat the error. Restricting to the first such episode per player-game leaves 29 later submissions and two repetitions. Opus's 60 refusal-path error decisions never receive a corrective request, and Opus accounts for 54 of the 63 repetitions in the full error population. With feedback this unequal, the aggregate cannot show that one configuration learns less than another. Within-auction ceilings are also revised: at G20 d40--46 Grok tightens its stated limit from \$235 to \$230, bids \$232 while acknowledging the overrun, then withdraws. None of these sequences shows durable learning. Testing that would need controlled feedback and comparable later opportunities.

\FloatBarrier
